\chapter{Aversion au risque}

La théorie de l'utilité de von Neumann et Morgenstern nous dit que
l'utilité associée à une loterie donnant un gain monétaire $x$
avec probabilité $\frac12$ et $y$ avec probabilité $\frac12$ est
la moyenne entre les utilités de $x$ et de $y$. Ce que la théorie
de von Neumann et Morgenstern ne nous dit pas, c'est comment se
comparent cette utilité avec celle du gain monétaire de $\frac12 x
+ \frac12 y$. C'est ce que nous étudions dans ce chapitre, qui
concerne les préférences de l'agent face au risque.

Les objectifs de ce chapitre sont :%
\begin{itemize}
\item D\'efinir un agent averse au risque ou qui aime le risque
\item Caractériser les fonctions d'utilité correspondantes. %
\item Étudier l'évolution du comportement face au risque avec la
richesse.
\item Mesurer cette aversion au risque ou cet amour du risque %
\end{itemize}

\section{Risque et utilités}

Les gains et les pertes sont mesurés de manière monétaire. On
considère donc des loteries à valeurs dans $\reals$, et soit $P=
\{$ loteries sur $\reals$ \`a support fini $\}$. En particulier,
pour $z\in \reals$, $\delta_z\in P$ représente la loterie qui
donne $z$ avec probabilité 1. Pour $z\in P$, $e(z)=\sum_x z(x).x$
est l'espérance de $p$. Pour $f\colon \reals\to\reals$, et $p\in
P$, $\E_p f$ est l'esp\'erance de $f$ sous $p$.

On se dote d'une relation de pr\'ef\'erences rationnelle $\succ$
sur $P$ repr\'esent\'ee par une fonction d'utilit\'e de von
Neumann--Morgenstern $u$.

On caractérise d'abord la croissance de l'utilité avec l'argent.
\begin{proposition}
$u$ est strictement croissante si et seulement si:
$$
\delta_z \succ \delta_{z'} \Leftrightarrow z > z'
$$
\end{proposition}

On supposera par la suite $u$ strictement croissante.

\begin{definition}~
\begin{itemize}
\item[$\bullet$] $\succ$ est {\bf averse au risque} si
$\delta_{e(p)}\succcurlyeq p$ pour tout $p\in P$.%
\item[$\bullet$] $\succ$ est {\bf strictement averse au risque} si
$\delta_{e(p)}\succ p$ pour tout $p\in P$.%
\item[$\bullet$] $\succ$ est {\bf neutre au risque} si
$\delta_{e(p)}\sim p$ pour tout $p\in P$.%
\item[$\bullet$] $\succ$ {\bf aime le risque} si
$\delta_{e(p)}\preccurlyeq p$ pour tout $p\in P$ %
\item[$\bullet$] $\succ$ {\bf aime strictement le risque} si
$\delta_{e(p)}\prec p$ pour tout $p\in P$.
\end{itemize}
\end{definition}
On a les caractérisations suivantes en termes de fonctions
d'utilité.
\begin{proposition}~
\begin{itemize}
\item[$\bullet$] $\succ$ est averse au risque si et seulement si
$u$ est
concave.%
\item[$\bullet$] $\succ$ est strictement averse au risque si et
seulement si
$u$ est strictement concave.%
\item[$\bullet$] $\succ$ est neutre au risque si $u$ est affine. %
\item[$\bullet$] $\succ$ aime le risque si $u$ est convexe. %
\item[$\bullet$] $\succ$ aime strictement le risque si $u$ est
strictement convexe.
\end{itemize}
\end{proposition}

\section{Risque et richesse}
Un agent aimera-t-il plus ou moins le risque lorsque sa richess
augmente ? Deux d\'efinitions sont possibles :
\begin{itemize}
\item Un m\^eme risque (en valeurs {\em absolues}) sera-t-il pris
par un agent plus riche, moins riche ?%
\item Un m\^eme risque (en valeurs {\em relatives}) sera-t-il pris
par un agent plus riche, moins riche ?
\end{itemize}

\subsection{Risque absolu}
On définit l'aversion au risque absolu en termes de préférences
avant de la caractériser en termes d'utilités.

\begin{definition}
$u$ (ou $\succ$) pr\'esente une {\em aversion au risque (absolu)
d\'ecroissante} [resp.\, croissante, resp.\, constante] si
$\forall$ $q\in P$, $z,w,w' \in \reals$ t.q.\, $w'>w$:
\begin{eqnarray*}
\E_{w+q}u>u(w+z) &\Rightarrow& \E_{w'+q}u>u(w'+z)\\
\mbox{resp.\,} &\Leftarrow& \\
\mbox{resp.\,} &\Leftrightarrow&
\end{eqnarray*}
\end{definition}

Dans le premier cas, si la loterie $q$ est pr\'ef\'er\'ee \`a la
somme $z$ s\^ure lorsque la richesse est $w$, elle l'est aussi si
la richesse est $w'>w$.


\begin{proposition} Supposons $u$ deux fois diff\'erentiable.
Alors $u$ pr\'esente une aversion au risque d\'ecroissante
[resp.\, croissante, resp.\, constante] si et seulement si la
fonction:
$$
\lambda = - \frac{u''}{u'}
$$
est d\'ecroissante [resp.\, d\'ecroissante, resp.\, constante]. \\
\end{proposition}

$\lambda$ s'appelle le {\bf coefficient d'aversion au risque de
Arrow-Pratt}. Il s'agit de la courbure de $u$ normalis\'ee par la
d\'eriv\'ee.


La proposition suivante caractérise l'aversion pour le risque
absolu constante:

\begin{proposition} Supposons $u$ deux fois diff\'erentiable. Alors
$u$ pr\'esente une aversion au risque constante si et seulement si
il existe $a>0$ et $b$ tels que:
$$
u(z)=
\begin{cases}
  az+b &\mbox{si }\;\; \lambda(z)\equiv0\cr
  -ae^{-\lambda z}+b  &\mbox{si }\;\; \lambda(z)\equiv\lambda>0
\end{cases}
$$
\end{proposition}

C'est-à-dire que les préférences sont représentables par une
fonction d'utilité $u$ qui peut être de deux types :
$$
u(z)=
\begin{cases}
  z &\mbox{si }\;\; \lambda(z)\equiv0\cr
  -e^{-\lambda z}  &\mbox{si }\;\; \lambda(z)\equiv\lambda>0
\end{cases}
$$

\subsection{Risque relatif}

\begin{definition}
$u$ (ou $\succ$) pr\'esente une {\em aversion au risque (relatif)
d\'ecroissante} [resp.\, croissante, resp.\, constante] si
$\forall$ $q\in P$ t.q.\ $q\geq 0$ avec pba.\, 1, $z>0$ $w'>w$:
\begin{eqnarray*}
\E_{wq}u>u(wz) &\Rightarrow& \E_{w'q}u>u(w'z)\\
\mbox{resp.\,} &\Leftarrow& \\
\mbox{resp.\,} &\Leftrightarrow&
\end{eqnarray*}
\end{definition}

Dans le premier cas, si on pr\'ef\`ere multiplier sa richesse par
la loterie $q$ plut\^ot que par $z$ lorsque la richesse est $w$,
on le pr\'ef\`ere aussi si la richesse est $w'>w$.

\begin{proposition} Supposons $u$ deux fois diff\'erentiable. Alors
$u$ pr\'esente une aversion au risque relatif d\'ecroissante
[resp.\, constante] si et seulement si la fonction:
$$
z\to  - \frac{zu''(z)}{u'(z)}
$$
est d\'ecroissante [resp.\, constante].
\end{proposition}

\begin{proposition} Supposons $u$ deux fois diff\'erentiable.
Alors $u$ pr\'esente une aversion au risque relatif constante si
et seulement si il existe $a>0$ et $b$ tels que:
$$
u(z)=
\begin{cases}
  a\ln(z)+b &\cr
  a \gamma z^{\gamma}+b  &\mbox{pour un certain }\gamma<1, \gamma\neq 0
\end{cases}
$$
\end{proposition}

C'est-à-dire que les préférences sont représentables par une
fonction d'utilité $u$ qui est d'un des trois types suivants :
$$
u(z)=
\begin{cases}
  \ln(z) & \cr
  - z^{\gamma}  &\mbox{pour un certain }\gamma<0\cr
    z^{\gamma}  &\mbox{pour un certain } 0<\gamma<1
\end{cases}
$$

\section{Application}
Supposons que vous ayez le choix entre les loteries suivantes:\\
\begin{minipage}{3.5cm}%
\begin{tabular}{p{0.4cm}p{0.4cm}p{0.4cm}}
&\rnode{a}{$\blue\bullet$}&\\[1.5cm]
\rnode{b}{$\blue\bullet$}&\rnode{c}{$\blue\bullet$}&\rnode{d}{$\blue\bullet$}\\[-0.4cm]
\red\tiny25\euro&\red\tiny150\euro&\red\tiny600\euro\\
\ncline[linewidth=1.3pt,linecolor=yellow]{a}{b}\aput[2pt]{0}(0.89){\tiny$0.7$}
\ncline[linewidth=1.3pt,linecolor=yellow]{a}{c}\aput[2pt]{0}(0.95){\tiny$0.2$}
\ncline[linewidth=1.3pt,linecolor=yellow]{a}{d}\aput[2pt]{0}(0.99){\tiny$0.1$}
\end{tabular}%
\end{minipage}%
\begin{minipage}{3.5cm}%
\begin{tabular}{p{0.4cm}p{0.4cm}p{0.4cm}}
&\rnode{a}{$\blue\bullet$}&\\[1.5cm]
\rnode{b}{$\blue\bullet$}&\rnode{c}{$\blue\bullet$}&\rnode{d}{$\blue\bullet$}\\[-0.4cm]
\red\tiny80\euro&\red\tiny90\euro&\red\tiny98\euro\\
\ncline[linewidth=1.3pt,linecolor=yellow]{a}{b}\aput[2pt]{0}(0.89){\tiny$0.2$}
\ncline[linewidth=1.3pt,linecolor=yellow]{a}{c}\aput[2pt]{0}(0.95){\tiny$0.58$}
\ncline[linewidth=1.3pt,linecolor=yellow]{a}{d}\aput[2pt]{0}(0.99){\tiny$0.22$}
\end{tabular}%
\end{minipage}%
\begin{minipage}{3.5cm}%
\begin{tabular}{p{0.4cm}p{0.4cm}p{0.4cm}}
&\rnode{a}{$\blue\bullet$}&\\[1.5cm]
\rnode{b}{$\blue\bullet$}&\rnode{c}{}&\rnode{d}{$\blue\bullet$}\\[-0.4cm]
\red\tiny2\euro&\red\tiny&\red\tiny105\euro\\
\ncline[linewidth=1.3pt,linecolor=yellow]{a}{b}\aput[2pt]{0}(0.89){\tiny$0.05$}
%\ncline[linewidth=1.3pt,linecolor=yellow]{a}{c}\aput[2pt]{0}(0.95){\tiny$0.58$}
\ncline[linewidth=1.3pt,linecolor=yellow]{a}{d}\aput[2pt]{0}(0.99){\tiny$0.95$}
\end{tabular}%
\end{minipage}\\
A-t-on une mani\`ere objective de choisir entre ces loteries~?

Une m\'ethode possible est de supposer une fonction d'utilité de
von Neumann et Morgenstern, avec une aversion au risque constante
{\em pour ces valeurs mon\'etaires}. Il reste \`a déterminer le
coefficient d'aversion au risque $\lambda$.

On peut déterminer $\lambda$ en répondant à la question suivante :
Combien donneriez-vous pour une loterie offrant $0$ avec
probabilit\'e $\frac12$ et $500$\,\euro\ avec probabilit\'e
$\frac12$ ?

Il est ensuite possible d'évaluer l'utilité pour chacune des
loteries, et de déterminer celle qui maximise l'utilité.

\chapter{Théorie de la décision dynamique}

\section{Introduction}
Nous considérons des situations dans lesquelles une suite de
décisions doit \^etre prise par l'agent. On montre comment décrire
ces situations par des arbres de décision, et comment décrire un
comportement de l'agent par une stratégie. On verra aussi comment
rechercher les stratégies optimales de l'agent, c'est-à-dire
celles qui maximisent son utilité.

\section{Décision dynamique en information parfaite}

\subsection{Description d'un problème de décision}
Un probl\`eme de décision dynamique en information parfaite est
donné par
\begin{itemize}
\item $T$ ensemble fini de n{\oe}uds, $Z\subset T$ n{\oe}uds
terminaux, et pour $t\not\in Z$ :
\begin{itemize}
\item  $i(t)\in \{N,J\}$ qui détermine si la nature ($N$) ou le
joueur ($J$) joue en $t$ ;%
\item L'ensemble de choix possibles au n{\oe}ud $t$, $A(t)$ ; %
\item Le n{\oe}ud successeur r\'esultant du choix $a$, $N(t,a)$ ;
\item Une distribution $D(t)$ sur $A(t)$ si $i(t)= N$.
\end{itemize}
\item $u:Z\to\reals$ fonctions d'utilit\'e du joueur ;
\end{itemize}

On définit ainsi les successeurs d'un n{\oe}ud
\begin{itemize}
\item $s(t)=\{N(t,a),a\in A(t)\}$ ensemble des successeurs directs
de $t$. ($\emptyset$ si $t\in Z$).%
\item $S(t)= s(t)\cup s(s(t)) \cup s(s(s(t)))...$ tous les
successeurs de $t$.
\end{itemize}

On suppose que le problème a une structure d'arbre, c'est à dire :

\begin{itemize}
\item Il n'y a pas de cycle : $\forall t, t\not\in S(t)$ %
\item Il existe un  n{\oe}ud initial $\exists\tilde t, S(\tilde
t)\cup\{\tilde t\}=T$ %
\item Chaque n{\oe}ud a un unique prédécesseur
\end{itemize}

\begin{exemple}
Exemple : casino.\\
$\bullet$ 1\`ere \'etape : Le joueur peut d\'ecider de jouer ou non.\\
S'il ne joue pas le jeu s'arr\^ete.\\
Sinon, la nature tire \`a pile ou face s'il gagne 1 ou perd 2.\\[0.5cm]
$\bullet$ 2\`eme \'etape : Le joueur peut d\'ecider de jouer ou non.\\
S'il ne joue pas le jeu s'arr\^ete.\\
S'il joue, la nature tire \`a pile ou face s'il gagne 3 ou perd 2
(en plus des gains et pertes de la première étape).

On représente ce jeu par l'arbre suivant :\\[0.5cm]
\begin{tabular}{p{1cm}p{1cm}p{1cm}p{1cm}p{1cm}p{1cm}}
&&&\rnode{a}{$\bullet$}&\rnode{c}{$0$}\\[0.7cm]
&&\rnode{b}{$\blue\diamond$}&&&\\[0.7cm]
&\rnode{d}{$\bullet$}&&\rnode{e}{$\bullet$}&\\[0.7cm]
&\rnode{f}{$\blue\diamond$}&\rnode{g}{\small$+1$}&\rnode{h}{$\blue\diamond$}&\rnode{i}{\small$-2$}&\\[0.7cm]
\rnode{j}{\small$+4$}&\rnode{k}{\small$-1$}&\rnode{l}{\small$+1$}&\rnode{m}{\small$-4$}\\[0.7cm]
\ncline[linewidth=0.7pt,linecolor=yellow]{-}{a}{b}\bput[1pt]{0}(0.5){\small$J_1$}%
\ncline[linewidth=0.7pt,linecolor=yellow]{-}{a}{c}\aput[1pt]{0}(0.5){\small$S_1$}%
\ncline[linewidth=0.7pt,linecolor=yellow]{-}{b}{d}\bput[1pt]{0}(0.5){\small$\frac12$}%
\ncline[linewidth=0.7pt,linecolor=yellow]{-}{b}{e}\aput[1pt]{0}(0.5){\small$\frac12$}%
\ncline[linewidth=0.7pt,linecolor=yellow]{-}{d}{f}\bput[1pt]{0}(0.5){\small$J_2$}%
\ncline[linewidth=0.7pt,linecolor=yellow]{-}{d}{g}\aput[1pt]{0}(0.5){\small$S_2$}%
\ncline[linewidth=0.7pt,linecolor=yellow]{-}{e}{h}\bput[1pt]{0}(0.5){\small$J_3$}%
\ncline[linewidth=0.7pt,linecolor=yellow]{-}{e}{i}\aput[1pt]{0}(0.5){\small$S_3$}%
\ncline[linewidth=0.7pt,linecolor=yellow]{-}{f}{j}\bput[1pt]{0}(0.5){\small$\frac12$}%
\ncline[linewidth=0.7pt,linecolor=yellow]{-}{f}{k}\aput[1pt]{0}(0.5){\small$\frac12$}%
\ncline[linewidth=0.7pt,linecolor=yellow]{-}{h}{l}\bput[1pt]{0}(0.5){\small$\frac12$}%
\ncline[linewidth=0.7pt,linecolor=yellow]{-}{h}{m}\aput[1pt]{0}(0.5){\small$\frac12$}%
%&\circlenode*{b}{\makebox[0.5cm]{{\red$B^{\blue11}$}}}&\circlenode*{e}{\makebox[0.5cm]{{\red$E^{\blue4}$}}}&&\\[0.1cm]
%&&&\circlenode*{h}{\makebox[0.5cm]{{\red$H^{\blue3}$}}}&\\[0.1cm]
%\circlenode*{a}{\makebox[0.5cm]{{\red$A^{\blue11}$}}}&\circlenode*{c}{\makebox[0.5cm]{{\red$C^{\blue7}$}}}%
\end{tabular}\\
Les n{\oe}uds $\bullet$ sont ceux où le joueur prend une décision.
Les décisions possibles au premier  n{\oe}ud de l'agent sont $J_1$
et $S_1$, $J_2$ et $S_2$ pour le deuxième, et $J_3$ et $S_3$ pour
le troisième. Les n{\oe}uds $\diamond$ sont ceux où la nature tire
au hasard. Les branches suivant ces n{\oe}uds sont affectées des
probabilités que la nature accorde à chacun d'entre eux. Les
n{\oe}uds affectés de nombres sont les n{\oe}uds terminaux, où
nous avons représenté le gain monétaire final de l'agent.

\end{exemple}

\subsection{Description d'un comportement : stratégie}

Une stratégie $s$ pour l'agent associe une décision dans $A(t)$ à
tout n{\oe}ud $t$ tel que $i(t)=J$. C'est donc une application de
$i^{-1}(J)$ vers $\cup_t A(t)$ telle que $s(t)\in A(t)$ pour tout
$t$.

A toute stratégie du joueur correspond une probabilité $p(s)$ sur
les n{\oe}uds terminaux, éléments de $Z$. Pour calculer $p(s)$, on
part du n{\oe}ud initial, et on descend progressivement d'une
étape à la fois, pour calculer les probabilités des successeurs,
en utilisant à un n{\oe}ud $t$ tel que $i(t)=N$ la probabilité
$D(t)$, et à un n{\oe}ud $t$ tel que $i(t)=J$ la décision $s(t)$.

On a donc une utilité espérée $U(s) = \E_{p(s)} u(z)$ pour chaque
stratégie. Une stratégie maximisatrice est une stratégie qui
maximise ce paiement espéré.

\begin{exemple}
On a dans le problème précédent 3 n{\oe}uds de décision pour
l'agent, avec à chacun d'entre eux 2 décisions possibles. Par
conséquent il y a 8 stratégies possibles. $J_1S_2S_3$ joue à la
première étape et s'arr\^ete ensuite. $J_1S_2J_3$ joue a la
première étape, puis continue si elle a perdu à la première, et
arr\^ete sinon. $S_1J_2J_3$ peut s'entendre comme ``ne pas jouer,
mais si on avait joué à la première étape, on aurait décidé de
continuer à la seconde''. Ce n'est pas la m\^eme chose que
$S_1S_2S_3$ qui d'arr\^ete à la première étape, mais décide aussi
de s'arr\^eter si on lui demande ce qu'elle ferait à la deuxième.

Les gains espérés de ces stratégies sont :
\begin{itemize}
\item $J_1S_2S_3$ donne $-\frac12$ %
\item $J_1S_2J_3$ donne $0$%
\item $S_1J_2J_3$ et $S_1S_2S_3$ donnent $0$.
\end{itemize}
\end{exemple}
%
\subsection{Principe d'induction à rebours}

Le principe d'induction à rebours permet de trouver une stratégie
maximisatrice. On part du bas de l'arbre, et pour les dernières
décisions de l'agent, on choisit la meilleure action. On remplace
ensuite le n{\oe}ud correspondant par l'utilité espérée étant
donnée cette meilleure action. On obtient ainsi un nouveau
problème de décision dynamique, auquel on applique de nouveau le
principe, et ainsi de suite... A la fin de la procédure, on aura
choisi une décision pour chaque n{\oe}ud, on aura donc obtenu une
stratégie dans le problème de décision. On a le résultat suivant :

\begin{proposition}
Toute stratégie obtenue par le principe d'induction à rebours est
maximisatrice.
\end{proposition}

\begin{exemple}
Question : Comment maximiser la probabilit\'e de finir avec un
gain d'au moins 1 dans l'exemple précédent ? Ceci revient à
chercher une stratégie maximisatrice lorsque l'utilité vaut 1 pour
les gains de 1 au moins, et 0 sinon. On a alors l'arbre :\\[0.5cm]
\begin{tabular}{p{1cm}p{1cm}p{1cm}p{1cm}p{1cm}p{1cm}}
&&&\rnode{a}{$\bullet$}&\rnode{c}{$0$}\\[0.7cm]
&&\rnode{b}{$\blue\diamond$}&&&\\[0.7cm]
&\rnode{d}{$\bullet$}&&\rnode{e}{$\bullet$}&\\[0.7cm]
&\rnode{f}{$\blue\diamond$}&\rnode{g}{\small$+1$}&\rnode{h}{$\blue\diamond$}&\rnode{i}{\small$0$}&\\[0.7cm]
\rnode{j}{\small$1$}&\rnode{k}{\small$0$}&\rnode{l}{\small$1$}&\rnode{m}{\small$0$}\\[0.7cm]
\ncline[linewidth=0.7pt,linecolor=yellow]{-}{a}{b}\bput[1pt]{0}(0.5){\small$J_1$}%
\ncline[linewidth=0.7pt,linecolor=yellow]{-}{a}{c}\aput[1pt]{0}(0.5){\small$S_1$}%
\ncline[linewidth=0.7pt,linecolor=yellow]{-}{b}{d}\bput[1pt]{0}(0.5){\small$\frac12$}%
\ncline[linewidth=0.7pt,linecolor=yellow]{-}{b}{e}\aput[1pt]{0}(0.5){\small$\frac12$}%
\ncline[linewidth=0.7pt,linecolor=yellow]{-}{d}{f}\bput[1pt]{0}(0.5){\small$J_2$}%
\ncline[linewidth=0.7pt,linecolor=yellow]{-}{d}{g}\aput[1pt]{0}(0.5){\small$S_2$}%
\ncline[linewidth=0.7pt,linecolor=yellow]{-}{e}{h}\bput[1pt]{0}(0.5){\small$J_3$}%
\ncline[linewidth=0.7pt,linecolor=yellow]{-}{e}{i}\aput[1pt]{0}(0.5){\small$S_3$}%
\ncline[linewidth=0.7pt,linecolor=yellow]{-}{f}{j}\bput[1pt]{0}(0.5){\small$\frac12$}%
\ncline[linewidth=0.7pt,linecolor=yellow]{-}{f}{k}\aput[1pt]{0}(0.5){\small$\frac12$}%
\ncline[linewidth=0.7pt,linecolor=yellow]{-}{h}{l}\bput[1pt]{0}(0.5){\small$\frac12$}%
\ncline[linewidth=0.7pt,linecolor=yellow]{-}{h}{m}\aput[1pt]{0}(0.5){\small$\frac12$}%
%&\circlenode*{b}{\makebox[0.5cm]{{\red$B^{\blue11}$}}}&\circlenode*{e}{\makebox[0.5cm]{{\red$E^{\blue4}$}}}&&\\[0.1cm]
%&&&\circlenode*{h}{\makebox[0.5cm]{{\red$H^{\blue3}$}}}&\\[0.1cm]
%\circlenode*{a}{\makebox[0.5cm]{{\red$A^{\blue11}$}}}&\circlenode*{c}{\makebox[0.5cm]{{\red$C^{\blue7}$}}}%
\end{tabular}

Appliquons le principe d'induction à rebours à cette arbre. On
voit qu'il faut décider $S_2$ au deuxième n{\oe}ud, et $J_3$ au
troisième. On obtient l'arbre dans lequel on a remplacé les
derniers n{\oe}uds de décisions par le gain espéré maximal :\\

\begin{tabular}{p{1cm}p{1cm}p{1cm}p{1cm}p{1cm}p{1cm}}
&&&\rnode{a}{$\bullet$}&\rnode{c}{$0$}\\[0.7cm]
&&\rnode{b}{$\blue\diamond$}&&&\\[0.7cm]
&\rnode{d}{$1$}&&\rnode{e}{$\frac12$}&\\[0.7cm]
\ncline[linewidth=0.7pt,linecolor=yellow]{-}{a}{b}\bput[1pt]{0}(0.5){\small$J_1$}%
\ncline[linewidth=0.7pt,linecolor=yellow]{-}{a}{c}\aput[1pt]{0}(0.5){\small$S_2$}%
\ncline[linewidth=0.7pt,linecolor=yellow]{-}{b}{d}\bput[1pt]{0}(0.5){\small$\frac12$}%
\ncline[linewidth=0.7pt,linecolor=yellow]{-}{b}{e}\aput[1pt]{0}(0.5){\small$\frac12$}%
%&\circlenode*{b}{\makebox[0.5cm]{{\red$B^{\blue11}$}}}&\circlenode*{e}{\makebox[0.5cm]{{\red$E^{\blue4}$}}}&&\\[0.1cm]
%&&&\circlenode*{h}{\makebox[0.5cm]{{\red$H^{\blue3}$}}}&\\[0.1cm]
%\circlenode*{a}{\makebox[0.5cm]{{\red$A^{\blue11}$}}}&\circlenode*{c}{\makebox[0.5cm]{{\red$C^{\blue7}$}}}%
\end{tabular}

On voit que dans ce jeu la meilleure décision est de jouer $J_1$
au premier n{\oe}ud. On a donc obtenu comme stratégie la stratégie
$J_1S_2J_3$ qui maximise la probabilité de terminer le jeu avec un
paiement d'au moins~1.
\end{exemple}

Une remarque importante sur la notion de stratégie et sur le
principe d'induction à rebours. Le principe d'action à rebours
indique quelle choix prendre, même aux n{\oe}uds qui ne devraient
pas être atteints. Par exemple, regardons le jeu suivant~:\\[0.5cm]

\begin{tabular}{p{1cm}p{1cm}p{1cm}p{1cm}p{1cm}p{1cm}}
&&&\rnode{a}{$\bullet$}&\rnode{c}{$0$}\\[0.7cm]
&&\rnode{b}{$\bullet$}&&&\\[0.7cm]
&\rnode{d}{$-1$}&&\rnode{e}{$-2$}&\\[0.7cm]
\ncline[linewidth=0.7pt,linecolor=yellow]{-}{a}{b}\bput[1pt]{0}(0.5){\small$J_1$}%
\ncline[linewidth=0.7pt,linecolor=yellow]{-}{a}{c}\aput[1pt]{0}(0.5){\small$S_1$}%
\ncline[linewidth=0.7pt,linecolor=yellow]{-}{b}{d}\bput[1pt]{0}(0.5){\small$J_2$}%
\ncline[linewidth=0.7pt,linecolor=yellow]{-}{b}{e}\aput[1pt]{0}(0.5){\small$S_2$}%
%&\circlenode*{b}{\makebox[0.5cm]{{\red$B^{\blue11}$}}}&\circlenode*{e}{\makebox[0.5cm]{{\red$E^{\blue4}$}}}&&\\[0.1cm]
%&&&\circlenode*{h}{\makebox[0.5cm]{{\red$H^{\blue3}$}}}&\\[0.1cm]
%\circlenode*{a}{\makebox[0.5cm]{{\red$A^{\blue11}$}}}&\circlenode*{c}{\makebox[0.5cm]{{\red$C^{\blue7}$}}}%
\end{tabular}

Le principe d'induction à rebours nous dit tout d'abord qu'il faut
choisir $J_2$ au deuxième n{\oe}ud, et transforme le problème en :\\[0.5cm]

\begin{tabular}{p{1cm}p{1cm}p{1cm}p{1cm}p{1cm}p{1cm}}
&&&\rnode{a}{$\bullet$}&\rnode{c}{$0$}\\[0.7cm]
&&\rnode{b}{$-1$}&&&\\[0.7cm]
\ncline[linewidth=0.7pt,linecolor=yellow]{-}{a}{b}\bput[1pt]{0}(0.5){\small$J$}%
\ncline[linewidth=0.7pt,linecolor=yellow]{-}{a}{c}\aput[1pt]{0}(0.5){\small$S$}%
%&\circlenode*{b}{\makebox[0.5cm]{{\red$B^{\blue11}$}}}&\circlenode*{e}{\makebox[0.5cm]{{\red$E^{\blue4}$}}}&&\\[0.1cm]
%&&&\circlenode*{h}{\makebox[0.5cm]{{\red$H^{\blue3}$}}}&\\[0.1cm]
%\circlenode*{a}{\makebox[0.5cm]{{\red$A^{\blue11}$}}}&\circlenode*{c}{\makebox[0.5cm]{{\red$C^{\blue7}$}}}%
\end{tabular}

On voit donc qu'il faut choisir $S_1$ au premier n{\oe}ud. La
procédure d'induction à rebours aboutit donc à la stratégie
$S_1J_2$. C'est parce qu'on a conclut qu'il fallait jouer $J_2$ et
que le gain correspondant est de $-1$ qu'on a pu ensuite conclure
qu'il fallait jouer $S_1$.

\section{Décision dynamique en information imparfaite}

On suppose maintenant que l'agent n'est pas complètement informé
de tous les coups de la nature. Au moment où il effectue ses
choix, il n'a qu'une information imparfaite sur sa position dans
l'arbre de décision. Nous commençons par décrire son information
dans l'arbre à l'aide de la notion d'ensembles d'information.

Un ensemble d'information est une classe d'équivalence sur les
n{\oe}uds $t$ tels que $i(t)=J$. Soit $h(t)$ la classe
d'équivalence du n{\oe}ud $t$, c'est-à-dire l'ensemble
d'information comprenant le n{\oe}ud $t$. Comme $h$ décrit des
relations d'équivalence on a~:
$$
t' \in h(t) \implies t \in h(t')
$$
 On suppose que l'agent conna\^it son ensemble de décisions à chaque n{\oe}ud.
L'ensemble d'information $h(t)$. On a la restriction suivante sur
la fonction qui à $t$ associe $h(t)$ :
$$
t'\in h(t) \implies A(t')=A(t)
$$

\begin{exemple}
Reprenons l'exemple du casino en supposant qu'après la première
étape, le joueur ne sait pas s'il a gagné ou non. On a l'arbre de
décision suivant :

\begin{tabular}{p{1cm}p{1cm}p{1cm}p{1cm}p{1cm}p{1cm}}
&&&\rnode{a}{$\bullet$}&\rnode{c}{$0$}\\[0.7cm]
&&\rnode{b}{$\blue\diamond$}&&&\\[0.7cm]
&\rnode{d}{$\bullet$}&&\rnode{e}{$\bullet$}&\\[0.7cm]
&\rnode{f}{$\blue\diamond$}&\rnode{g}{\small$+1$}&\rnode{h}{$\blue\diamond$}&\rnode{i}{\small$-2$}&\\[0.7cm]
\rnode{j}{\small$+4$}&\rnode{k}{\small$-1$}&\rnode{l}{\small$+1$}&\rnode{m}{\small$-4$}\\[0.7cm]
\ncline[linewidth=0.7pt,linecolor=yellow]{-}{a}{b}\bput[1pt]{0}(0.5){\small$J$}%
\ncline[linewidth=0.7pt,linecolor=yellow]{-}{a}{c}\aput[1pt]{0}(0.5){\small$S$}%
\ncline[linewidth=0.7pt,linecolor=yellow]{-}{b}{d}\bput[1pt]{0}(0.5){\small$\frac12$}%
\ncline[linewidth=0.7pt,linecolor=yellow]{-}{b}{e}\aput[1pt]{0}(0.5){\small$\frac12$}%
\ncline[linewidth=0.7pt,linecolor=yellow]{-}{d}{f}\bput[1pt]{0}(0.5){\small$J$}%
\ncline[linewidth=0.7pt,linecolor=yellow]{-}{d}{g}\aput[1pt]{0}(0.5){\small$S$}%
\ncline[linewidth=0.7pt,linecolor=yellow]{-}{e}{h}\bput[1pt]{0}(0.5){\small$J$}%
\ncline[linewidth=0.7pt,linecolor=yellow]{-}{e}{i}\aput[1pt]{0}(0.5){\small$S$}%
\ncline[linewidth=0.7pt,linecolor=yellow]{-}{f}{j}\bput[1pt]{0}(0.5){\small$\frac12$}%
\ncline[linewidth=0.7pt,linecolor=yellow]{-}{f}{k}\aput[1pt]{0}(0.5){\small$\frac12$}%
\ncline[linewidth=0.7pt,linecolor=yellow]{-}{h}{l}\bput[1pt]{0}(0.5){\small$\frac12$}%
\ncline[linewidth=0.7pt,linecolor=yellow]{-}{h}{m}\aput[1pt]{0}(0.5){\small$\frac12$}%
\psellipse[linecolor=green](2.9,3.66)(1.8,0.3)
%&\circlenode*{b}{\makebox[0.5cm]{{\red$B^{\blue11}$}}}&\circlenode*{e}{\makebox[0.5cm]{{\red$E^{\blue4}$}}}&&\\[0.1cm]
%&&&\circlenode*{h}{\makebox[0.5cm]{{\red$H^{\blue3}$}}}&\\[0.1cm]
%\circlenode*{a}{\makebox[0.5cm]{{\red$A^{\blue11}$}}}&\circlenode*{c}{\makebox[0.5cm]{{\red$C^{\blue7}$}}}%
\end{tabular}

Un premier ensemble d'information est donné par le n{\oe}ud au
début du jeu. Un deuxième contient les deux n{\oe}uds qui suivent
le choix de la nature après une première décision de jouer.
\end{exemple}

Une stratégie $s$ est une fonction de l'ensemble des {\em
ensembles d'information} vers les décisions possibles. On peut
voir $s$ comme une application de $i^{-1}(J)$ vers $\cup_t A(t)$
telle que :%
\begin{itemize}
\item $s(t)= s(t')$ si $h(t) = h(t')$ ($s(t)$ ne dépend que de
l'ensemble d'information de $t$)%
\item $s(t)\in A(t)$ pour tout $t$.
\end{itemize}

\begin{exemple}
On a deux ensembles d'information, et deux décisions possibles
pour chacun d'entre eux, soient 4 stratégies. $JS$, $JJ$, $SJ$ et
$SS$. Les paiement espérés correspondants sont :
\begin{itemize}
\item $JS$ donne $-\frac12$%
\item $JJ$ donne $0$%
\item $SJ$ et $SS$ donnent $0$.
\end{itemize}
\end{exemple}

\begin{exemple}
Rechercher les stratégies qui maximisent la probabilité de gagner
au moins 1 dans ce dernier jeu.
\end{exemple}
